\paragraph{Text-only datasets.} Our evaluation datasets are text-only by design (Section~\ref{sec:experimental_setup}), so that diversity comes from domain and task rather than modality and so that small models remain competitive under the parameter budget. This isolates domain generalization but leaves open how the approach transfers to image, video, audio, and multimodal datasets. Nothing in the pipeline is text-specific -- the dataset card is summarized and embedded regardless of modality, and ArtifactLinker's graph already contains vision and speech artifacts -- but the neighborhoods are sparser there, evaluation metrics are more heterogeneous, and the Evaluation Agent would need to handle modality-specific preprocessing. Whether the relative strengths of cosine retrieval and direct GNN scoring hold in these regions of the graph is an open question.

\paragraph{Closed candidate pool.} Both retrieval approaches, and the coding-agent baseline, choose from the 9{,}848 models present in ArtifactLinker's graph (5{,}830 of which satisfy the 14B-parameter constraint used in our runs). This closed list is what makes the comparison controlled and the GNN applicable, but it excludes models released after the graph snapshot and models that were never evaluated on a Hugging Face dataset, which are often the ones a practitioner most wants to know about. Lifting this restriction is not simply a matter of removing the list: with an unconstrained pool, every method degenerates to recommending the strongest available model, so meaningful comparison still requires practical constraints (parameter budget, license, inference cost) that keep the choice non-trivial. Extending the graph incrementally with newly released models, and evaluating the system on an open pool under such constraints, is a natural next step.

\paragraph{Cold-start for the graph, not for the LLM.} We define a cold-start dataset as one absent from ArtifactLinker's graph, but the datasets we evaluate on are public Hugging Face benchmarks, and the LLMs in the loop -- GPT-4o for summarization, GPT-5.5 for recommendation, and the coding-agent baseline -- have almost certainly seen these datasets, their leaderboards, and papers evaluating models on them during pretraining. Part of what the recommender does may therefore be recall rather than reasoning over the retrieved evidence, which inflates every LLM-in-the-loop condition relative to the raw-GNN baseline. What we measure is generalization to datasets \emph{outside the graph}, not to datasets outside the LLM's training distribution; the setting a practitioner most cares about -- private or newly collected data -- remains untested. A held-out set of datasets released after the LLMs' training cutoffs would separate the two, and we leave this to future work.

\paragraph{Baseline confounds evidence with model.} Our coding-agent baseline runs Claude Opus~4.7, whereas the retrieval conditions use GPT-5.5. A gap between the baseline and the merged condition therefore cannot be attributed to graph evidence alone, since it may partly reflect differences between the two models. The clean ablation is the same LLM, prompt, and candidate list with the evidence blocks removed; the baseline should be read as a strong practical reference point -- what a practitioner gets today by asking a coding agent -- rather than as a controlled measurement of the value of graph evidence.
